% -*- coding: utf-8; -*-

\chapter{Introdução}

O desenvolvimento de projetos de produção de petróleo e gás natural requer o planejamento estruturado de diversas etapas ao longo do ciclo de vida do campo, abrangendo exploração, avaliação, desenvolvimento, operação e abandono. Na fase exploratória, o principal objetivo é identificar acumulações potenciais de hidrocarbonetos por meio de levantamentos geofísicos, como sísmica, gravimetria e análises geológicas. Caso sejam encontrados indícios de petróleo e gás, inicia-se a fase de avaliação, destinada a determinar a viabilidade técnica e econômica da produção. Uma vez comprovado o potencial do reservatório, passa-se à fase de desenvolvimento, na qual são elaborados planos para instalação de plataformas, perfuração e completação de poços, buscando maximizar a recuperação de hidrocarbonetos com o menor custo possível.

Campanhas de perfuração otimamente planejadas são fundamentais para aumentar a recuperação de óleo e gerar valor ao projeto. Elas ocorrem em diferentes estágios do ciclo de vida do campo, desde a exploração até o desenvolvimento e perfuração complementar. Uma parcela significativa do risco econômico está associada à incerteza geológica, razão pela qual as equipes de subsuperfície são frequentemente responsáveis por otimizar parâmetros da campanha de perfuração, como o número de poços, suas localizações e a sequência de perfuração.

As empresas de petróleo utilizam as informações obtidas dos poços perfurados para ajustar as etapas subsequentes da campanha, incorporando o aprendizado adquirido. No entanto, métodos clássicos de otimização geralmente não consideram explicitamente esse processo de aprendizado sequencial e flexibilidade temporal. Assim, considerar a aquisição sequencial de informações geológicas e a tomada de decisão adaptativa durante as campanhas de perfuração pode agregar valor significativo ao projeto, mesmo sob condições de incerteza na caracterização geológica e nos parâmetros econômicos advindos do mercado e que influenciam diretamente o valor do projeto. O desempenho insatisfatório de muitos empreendimentos petrolíferos decorre da baixa aplicação de recursos na tomada de decisão e da ausência de investimentos adequados para estimar o retorno esperado do projeto, o que resulta na superestimação dos ganhos ou na subestimação das perdas.

Na literatura há consenso que três tipos de riscos devem ser considerados: perda de oportunidade (quando um prospecto é incorretamente considerado inviável), desenvolvimento não comercial (quando um campo é desenvolvido mesmo sendo antieconômico) e desenvolvimento subótimo (quando o campo é desenvolvido com retorno inferior ao máximo possível). A mitigação desses riscos depende da obtenção de informações adicionais que reduzam as incertezas relevantes. Uma metodologia consolidada para avaliar a aquisição de informação é a Análise de Valor da Informação, que relaciona a incerteza do reservatório às consequências econômicas potenciais. Embora o conceito seja amplamente difundido, sua aplicação prática torna-se complexa quando há múltiplos atributos incertos e informações incompletas ou pouco confiáveis. A integração entre Valor da Informação, Teoria de Opções Reais e modelos estocásticos de preços constitui uma abordagem robusta para apoiar decisões sob incerteza no setor de Exploração e Produção. A Teoria de Opções Reais incorpora o valor da flexibilidade gerencial, como adiar, expandir, suspender ou abandonar projetos conforme novas informações são obtidas e a favorabilidade do mercado. Entre os métodos numéricos, o Least Squares Monte Carlo destaca-se por permitir a estimação do valor ótimo de decisões sequenciais, considerando múltiplas fontes de incerteza e caminhos simulados de preços e custos. As incertezas de mercado podem ser representadas por modelos estocásticos como o Movimento Geométrico Browniano ou pelo modelo de dois fatores de Schwartz e Smith, que decompõe o logaritmo do preço do petróleo em dois componentes: um de longo prazo (nível de equilíbrio) e outro de curto prazo (desvio transitório). Essa formulação captura tanto a reversão à média quanto a volatilidade transitória, permitindo análises econômicas mais realistas e robustas.

Nos últimos anos, no entanto, a transformação digital e o avanço da Inteligência Artificial vêm modificando profundamente o processo decisório na indústria de petróleo e gás. Essa transformação, frequentemente associada à chamada quarta revolução industrial, é impulsionada pela convergência entre tecnologias físicas e digitais, como Aprendizado de Máquina, robótica e sistemas autônomos. O Aprendizado de Máquina tem demonstrado grande potencial em aplicações de engenharia de reservatórios, permitindo aprimorar a análise de grandes volumes de dados e identificar padrões complexos. Em particular, o Aprendizado Profundo, por meio de redes neurais artificiais e redes neurais convolucionais, vem sendo amplamente utilizado para modelar relações não lineares entre as propriedades petrofísicas do reservartório e desempenho na produção dos poços. Tais técnicas têm se mostrado eficazes para estimar a produtividade de poços, permeabilidade efetiva e posicionamento ótimo sob incerteza geológica. Estudos recentes demonstram a integração de redes neurais convolucionais com esquemas de otimização robusta, cujo objetivo é determinar a melhor localização de poços produtores em reservatórios sob incerteza. Nessas aplicações, as redes neurais convolucionais são treinadas para correlacionar propriedades petrofísicas espaciais com a produção acumulada de hidrocarbonetos, fornecendo estimativas de produtividade em diferentes realizações equiprováveis.

Com base nesses avanços, esta proposta de tese propõe a utilização de Aprendizado de Máquina para estimar o Valor da Informação Sequencial e apoiar a delimitação de campos de petróleo tendo como parâmetro de decisão uma medida de redução da incerteza. Embora o uso de Aprendizado de Máquina na fase de desenvolvimento seja uma realidade, sua aplicação na etapa de delimitação ainda é recente e promissora. A integração de métodos de Inteligência Artificial com a análise econômica de projetos exploratórios representa uma nova fronteira de pesquisa, permitindo quantificar de forma mais precisa o impacto do aprendizado geológico e econômico sobre o valor total do projeto.

A motivação para o desenvolvimento desse trabalho ...

Os resultados preliminares são apresentados no capítulo ...

A principal contribuição desse trabalho é ...

\section{Objetivos}
Os objetivos desse trabalho são:
\begin{itemize}
    \item[\textbullet] Desenvolver um novo método para estimar a locação de poços exploratórios na fase de delimitação com incertezas geológicas utilizando Valor da Informação e Medida de Aprendizagem;
    \item[\textbullet] Desenvolver modelos de Deep Learning via Redes Neurais Convolucionais que levem em consideração o novo método e as incertezas geológicas;
    \item[\textbullet] Aplicar o método desenvolvido em modelos de reservatórios com características similares aos reservatórios reais.
\end{itemize}

\section{Publicações}
No início do doutorado o aluno apresentou o trabalho \textbf{Pareto optimization via genetic algorithm for selecting a portfolio of oil and gas exploration projects} no \textit{2023 International Conference on Algorithms, Computing and Data Processing} e no \textit{$18^{th}$ International Congress of the Brazilian Geophysical Society \& Expogef}. O trabalho foi resultado da disciplina Computação Evolucionária ministrada pela Profa. Manoela Kohler e Prof. Marco Aurélio: \url{https://ieeexplore.ieee.org/document/10483593}.

Um estudo orientado realizado com os professores Marco Antonio Dias e Marco Aurélio resultou no trabalho \textbf{Learning, Uncertainty Reduction and Optimizing the Location of Delimitation Wells in Hydrocarbon E\&P} que foi apresentado na \textit{$27^{th}$ Annual International Real Options Conference} em Bologna/Italy: \url{https://realoptions.org/programAbstracts2024/_proceedings.html}. 

O aluno também participou do Programa de Doutorado Sanduíche no Exterior (PDSE/CAPES), na Universidade de Stavanger na Noruega, sob a supervisão do Prof. Reidar Bratvold, uma referência mundial em Métodos de Apoio à Decisão com aplicações na indústria de petróleo. O produto desse estágio foi um artigo científico que atualmente está em fase final de ajuste para submissão. Esse mesmo trabalho será apresentado na \textit{Brazilian Geophysical Society Conference - Rio'25} com o título \textbf{Real Options Approach for Economic Evaluation of Oil and Gas E\&P Projects Under Geological and Market Uncertainties}: \url{https://rio25.sbgf.org.br/Pages/technicalProgram.php#start}.

\section{Organização da Proposta}
No capítulo 2 é apresentado a fundamentação teórica para o entendimento do trabalho desenvolvido. No capítulo 3 é apresentado a nova técnica para lidar com amostras ruidosas RDS, assim como os modelos RDS-C e RDS-J. No capítulo 4 é apresentado o novo modelo RDS-Contrastive. No capítulo 5 é apresentado os detalhes experimentais para avaliação dos modelos. Também é apresentado os resultados dos modelos nos experimentos propostos, além disso é realizado uma análise nas vantagens e desvantagens de cada modelo. Ao final do capítulo 5 é apresentado o Estudo de Caso em detalhes e os resultados preliminares obtidos. No capítulo 6 é apresentado a conclusão e os próximos passos a serem realizados para conclusão desta tese. 











